\documentclass[10pt,a4paper]{amsart}
\usepackage[margin=19mm]{geometry}
\usepackage{amsmath,amssymb,mathtools}
\usepackage[scr=rsfs,cal=euler]{mathalfa}
\usepackage{xfrac}
\usepackage[unicode,hidelinks,bookmarks=false]{hyperref}
\newcommand{\ie}{i.e.\ }
\newcommand{\cf}{cf.\ }
\newcommand{\resp}{respectively\ }
\newcommand{\Subsection}{\subsection}
\providecommand{\nobreakdash}{\nobreak\hbox{-}}
\setlength{\emergencystretch}{2em}
\multlinegap0pt
\usepackage{bm}
\usepackage{bbm}
\usepackage{dsfont}
\usepackage{xspace}
\usepackage{cancel}
\usepackage{mathtools}
\usepackage{amsthm}
\makeatletter
\@ifundefined{lemma}{\newtheorem{lemma}{Lemma}}{}
\@ifundefined{remark}{\newtheorem{remark}{Remark}}{}
\makeatother
\makeatletter
\expandafter\ifx\csname proof\endcsname\relax
\renewenvironment{proof}{\noindent{\bfseries Proof.}}{\hfill$\blacksquare$}
\fi
\makeatother
\title[Cumulant-based quantum relative R\'enyi functional]{Cumulant-based quantum relative R\'enyi functional}
\author{Atirat Meunson, Tanapat Deesuwan}
\date{Source: arXiv:2606.31205v1 (CC BY 4.0)}
\begin{document}
\maketitle

\noindent\emph{Source and licence.} Cumulant-based quantum relative R\'enyi functional. Version arXiv:2606.31205v1 (CC BY 4.0), distributed under the Creative Commons Attribution 4.0 International licence, as recorded in the arXiv licence metadata for this exact version and shown on the abstract page https://arxiv.org/abs/2606.31205v1. Full rights notice: licenses/arxiv-2606.31205-CC-BY-4.0.txt.\par\smallskip

\noindent\textit{Editorial scope.} Counted unit: the Remark whose source label is \texttt{None} in the article (source0.tex lines 821--825, source label \texttt{None}), delivered together with the article's own complete original proof of it (source0.tex lines 827--1025, one \texttt{proof} environment of 199 lines). No mathematics is written, repaired or re-derived: statement and proof are the article's own text line for line. Required context retained uncounted: Geneq (lines 648--650); Lem1 (lines 794--805). Every definition, display and label of those lines is retained unchanged. Omitted from this package and not redistributed: 1--647, 651--793, 806--820, 826--826, 1026--2516. Nothing inside a retained excerpt is omitted or altered: every retained body is delivered exactly as the article's lines are written, so no omission receipt of any kind accompanies this package. Citations: the retained excerpts carry no citation command at all, so no bibliography entry of the article is transcribed and no reference is redistributed. Reference adaptation: none was needed and none was made; every reference inside the delivered unit resolves inside it, so no unresolved reference or citation is produced. Printed slips kept literal: no printed word, symbol, hypothesis or number of the counted statement or of its proof was altered. Layout and class: the article's own class is \texttt{revtex4-2} with packages including babel, inputenc, todonotes, amsthm, mathtools, physics, xcolor, graphicx, geometry, adjustbox, placeins, fontenc, lipsum, csquotes; the wrapper is the lane's pinned \texttt{amsart} editorial wrapper at 10pt on A4, which restates the theorem-like environments the retained text needs and adds the article's own macro definitions that the retained text needs (none), with extra wrapper packages (bm, bbm, dsfont, xspace, cancel, mathtools). The delivered numbering is the unit's own. Assets: no retained span contains a figure environment, a graphics-inclusion command, a PDF-page inclusion, a table or a TikZ picture; no image file is copied into this package, the package declares no asset dependency and no image file is redistributed with it. Licence: version arXiv:2606.31205v1 of this article is distributed under the Creative Commons Attribution 4.0 International licence, as recorded in the arXiv licence metadata for this exact version and shown on the abstract page https://arxiv.org/abs/2606.31205v1. A full rights notice is delivered at licenses/arxiv-2606.31205-CC-BY-4.0.txt. Characters: every retained excerpt is pure ASCII, so the delivered engine, XeLaTeX with its default Latin Modern text font, reports no missing character.
\begin{equation}\label{Geneq}
\Phi_{\rho,\sigma}(\alpha-1):=\operatorname{Tr}\!\left[ \rho\,e^{(\alpha-1)(\ln\rho-\ln\sigma)}\right].
\end{equation}

\begin{lemma}[H\"{o}lder inequality for Schatten classes]\label{Lem1}
Let $A \in \mathcal{L}_p(\mathcal{H})$ and $B \in \mathcal{L}_q(\mathcal{H})$ denote 
the Schatten $p$- and $q$-class operators respectively, with $1/p + 1/q = 1$.
Then
\begin{equation}
    \big|\operatorname{Tr}(AB)\big| \le \|A\|_p \, \|B\|_q.
\end{equation}
In particular, for $p=1$ and $q=\infty$,
\begin{equation}
    \big|\operatorname{Tr}(AB)\big| \le \|A\|_1 \, \|B\|_\infty.
\end{equation}
\end{lemma}

\begin{remark}
The full-rank assumption is imposed only for the continuity argument presented here.
It guarantees that the spectra of the states remain uniformly bounded away from zero, so that the logarithm is continuous under functional calculus and the operator \(\ln\rho-\ln\sigma\)
depends continuously on the pair \((\rho,\sigma).\)
\end{remark}

\begin{proof}
Let us consider the generating functional 
$\Phi_{\rho,\sigma}(\alpha-1)$ in~\eqref{Geneq}. In finite-dimensional Hilbert spaces, the Schatten norms satisfy $\|A\|_\infty \le \|A\|_1$ for any operator $A$. Therefore, convergence in trace norm implies convergence in operator norm, which yields $\|\rho_n-\rho\|_\infty \to 0$ and $\|\sigma_n-\sigma\|_\infty \to 0$. 

Since $\rho$ and $\sigma$ are of full rank, their minimal eigenvalues are strictly positive.
Moreover, since density operators have spectrum contained in $[0,1]$, and full-rankness guarantees a strictly positive lower spectral bound, there exists $c>0$ and $N_0 \in \mathbb{N}$ such that for all $n \ge N_0$,
\begin{equation}
\rho_n,\sigma_n,\rho,\sigma \ge cI.
\end{equation}
By continuity of eigenvalues with respect to the operator norm, such a uniform bound is stable under perturbation. 

Thus, their spectra lie in the compact interval $[c,1]$. 
Since $\log$ is continuous on $[c,1]$, functional calculus yields 
$\|\ln\rho_n-\ln\rho\|_\infty \to 0$ and $\|\ln\sigma_n-\ln\sigma\|_\infty \to 0$.

Set
\begin{equation}
    X_n := (\alpha-1)(\ln\rho_n-\ln\sigma_n)
\end{equation}
and 
\begin{equation}
    X := (\alpha-1)(\ln\rho-\ln\sigma).
\end{equation}
Using linearity of the operator norm and the triangle inequality, we estimate 
$\|X_n - X\|_\infty$ as follows. First, substituting the definitions of $X_n$ and 
$X$ and factoring out $|\alpha-1|$ via $\|cA\|_\infty = |c|\|A\|_\infty$,
\begin{align}
    \|X_n-X\|_{\infty} 
    &= \bigl\|(\alpha-1)(\ln\rho_n-\ln\sigma_n) \nonumber \\
    &\quad - (\alpha-1)(\ln\rho-\ln\sigma)\bigr\|_{\infty} \nonumber \\
    &= |\alpha-1|\,\Bigl\|(\ln\rho_n - \ln\sigma_n) \nonumber \\
    &\quad - (\ln\rho - \ln\sigma)\Bigr\|_{\infty}.
\end{align}
Regrouping the terms inside the norm as 
$(\ln\rho_n - \ln\rho) - (\ln\sigma_n - \ln\sigma)$
and applying the triangle inequality $\|A - B\|_\infty \le \|A\|_\infty + \|B\|_\infty$,
which follows from $\|A+(-B)\|_\infty \le \|A\|_\infty + \|-B\|_\infty$ and 
homogeneity $\|-B\|_\infty = \|B\|_\infty$, we obtain
\begin{align}
    \|X_n-X\|_{\infty} 
    &\le |\alpha-1|\,\bigl(\|\ln\rho_n - \ln\rho\|_{\infty} \notag\\
    &\qquad + \|\ln\sigma_n - \ln\sigma\|_{\infty}\bigr).
\end{align}
Since $\|\ln\rho_n - \ln\rho\|_\infty \to 0$ and $\|\ln\sigma_n - \ln\sigma\|_\infty \to 0$ 
by functional calculus, and $|\alpha-1|$ is a finite constant, we conclude
\begin{equation}
    \|X_n - X\|_\infty \to 0.
\end{equation}
The exponential map is continuous with respect to the operator norm on bounded 
operators, that is, $\|A_n - A\|_\infty \to 0$ implies 
$\|e^{A_n} - e^{A}\|_\infty \to 0$. Hence,
\begin{equation}
    \|e^{X_n}-e^{X}\|_{\infty}\to 0.
\end{equation}
Since $X_n\to X$ in operator norm, by definition of convergence 
applied with $\varepsilon = 1$, there exists $N_1\in\mathbb{N}$ such that
\begin{equation}
    \|X_n - X\|_{\infty} < 1 \qquad \text{for all } n \ge N_1.
\end{equation}
Applying the triangle inequality,
\begin{align}
    \|X_n\|_{\infty} 
    &= \|X_n - X + X\|_{\infty} \notag\\ 
    &\le \underbrace{\|X_n - X\|_{\infty}}_{<\, 1} + \|X\|_{\infty} \notag\\
    &< \|X\|_{\infty} + 1
    \qquad \text{for all } n\ge N_1.
\end{align}
Set $r := \|X\|_{\infty}+1 > 0$, so that $\|X_n\|_{\infty} < r$ for all $n \ge N_1$.
To bound $\|e^{X_n}\|_\infty$, we use the power series expansion
\begin{equation}
    e^{X_n} = \sum_{k=0}^{\infty} \frac{X_n^k}{k!} 
\end{equation}
Taking the operator norm of both sides and applying the triangle inequality for 
infinite sums $\bigl\|\sum_{k=0}^{\infty} A_k\bigr\|_\infty \le \sum_{k=0}^{\infty} 
\|A_k\|_\infty$, we obtain
\begin{equation}
    \|e^{X_n}\|_{\infty} 
    = \left\|\sum_{k=0}^{\infty} \frac{X_n^k}{k!}\right\|_\infty
    \le \sum_{k=0}^{\infty} \frac{\|X_n^k\|_{\infty}}{k!}.
\end{equation}
Applying submultiplicativity of the operator norm, that is, 
$\|AB\|_\infty \le \|A\|_\infty\|B\|_\infty$, iteratively gives
\begin{equation}
\begin{aligned}
    \|X_n^k\|_\infty 
    &= \|X_n \cdot X_n \cdots X_n\|_\infty \\
    &\le \|X_n\|_\infty \cdot \|X_n\|_\infty \cdots \|X_n\|_\infty \\
    &= \|X_n\|_\infty^k.
\end{aligned}
\end{equation}
Substituting this bound,
\begin{equation}
    \sum_{k=0}^{\infty} \frac{\|X_n^k\|_{\infty}}{k!} 
    \le \sum_{k=0}^{\infty} \frac{\|X_n\|_{\infty}^k}{k!}.
\end{equation}
Since $\|X_n\|_\infty$ is a real number, the right-hand side is exactly the 
power series of the real exponential, so
\begin{equation}
    \sum_{k=0}^{\infty} \frac{\|X_n\|_{\infty}^k}{k!} = e^{\|X_n\|_{\infty}}.
\end{equation}
Combining the above,
\begin{equation}
    \|e^{X_n}\|_{\infty} \le \sum_{k=0}^{\infty} \frac{\|X_n^k\|_\infty}{k!}
    \le \sum_{k=0}^{\infty} \frac{\|X_n\|_{\infty}^k}{k!} 
    = e^{\|X_n\|_{\infty}}.
\end{equation}
Because $e^{(\cdot)}$ is increasing on $\mathbb{R}$ and $\|X_n\|_\infty < r$ 
for all $n \ge N_1$, we conclude
\begin{equation}
    \|e^{X_n}\|_{\infty} \le e^{\|X_n\|_{\infty}} \le e^{r}
    \qquad \text{for all } n\ge N_1.
\end{equation}

To establish convergence of $\Phi_{\rho_n,\sigma_n}(\alpha-1)$,
consider the difference
\begin{equation}
    \Delta_n =
    \big|
    \operatorname{Tr}(\rho_n e^{X_n})
    -
    \operatorname{Tr}(\rho e^{X})
    \big|.
\end{equation}
By adding and subtracting $\operatorname{Tr}(\rho e^{X_n})$ and applying the
triangle inequality,
\begin{equation}
    \Delta_n \le \big|\operatorname{Tr}((\rho_n-\rho)e^{X_n})\big|+
    \big|\operatorname{Tr}(\rho(e^{X_n}-e^{X}))\big|.
\end{equation}
From Lemma~\ref{Lem1}, we obtain 
\begin{equation}
    \big|\operatorname{Tr}((\rho_n-\rho)e^{X_n})\big|\le\|\rho_n-\rho\|_1 \, \|e^{X_n}\|_\infty,
\end{equation}
and, using $\|\rho\|_1 = \operatorname{Tr}(\rho) = 1$,
\begin{equation}
    \big|\operatorname{Tr}(\rho(e^{X_n}-e^{X}))\big|\le
    \|\rho\|_1 \, \|e^{X_n}-e^{X}\|_{\infty} = \|e^{X_n}-e^{X}\|_{\infty}.
\end{equation}
Thus, 
\begin{equation}
    \Delta_n\le \|\rho_n-\rho\|_1 \, \|e^{X_n}\|_{\infty}
    + \|e^{X_n}-e^{X}\|_{\infty}.
\end{equation}
Since $\|\rho_n-\rho\|_1 \to 0$, for any $\varepsilon_1>0$ there exists $N_2\in\mathbb{N}$ 
such that
\begin{equation}
    \|\rho_n-\rho\|_1 < \frac{\varepsilon_1}{e^{r}}
    \quad \text{for all } n\ge N_2.
\end{equation}
Consequently, for all $n\ge \max\{N_1,N_2\}$,
\begin{equation}
    \|\rho_n-\rho\|_1\,\|e^{X_n}\|_{\infty} < \varepsilon_1.
\end{equation}    

On the other hand, the continuity of the exponential map in operator norm implies
$\|e^{X_n}-e^{X}\|_{\infty}\to 0$. Hence, for any $\varepsilon_2>0$ there exists 
$N_3\in\mathbb{N}$ such that
\begin{equation}
    \|e^{X_n}-e^{X}\|_{\infty} < \varepsilon_2
    \quad \text{for all } n\ge N_3.
\end{equation}
Letting $N := \max\{N_0, N_1, N_2, N_3\}$, we obtain for all $n\ge N$,
\begin{equation}
    \Delta_n \le \|\rho_n-\rho\|_1\,\|e^{X_n}\|_{\infty}+
    \|e^{X_n}-e^{X}\|_{\infty}<\varepsilon_1+\varepsilon_2.
\end{equation}
Since $\varepsilon_1$ and $\varepsilon_2$ are arbitrary with
$\varepsilon_1+\varepsilon_2=\varepsilon$, it follows that
$\Delta_n\to 0$ as $n\to\infty$. Hence,
\[
    \Phi_{\rho_n,\sigma_n}(\alpha-1)
    \longrightarrow
    \Phi_{\rho,\sigma}(\alpha-1).
\]
Since $e^{X_n}\succ 0$ for all $n$ and $e^X\succ 0$, and both $\rho_n$ and $\rho$
are density operators, we have
\begin{equation}
    \Phi_{\rho_n,\sigma_n}(\alpha-1)=\operatorname{Tr}(\rho_n e^{X_n})>0
\end{equation}
and
\begin{equation}
     \Phi_{\rho,\sigma}(\alpha-1)=\operatorname{Tr}(\rho e^{X})>0
\end{equation}
for all $n\in\mathbb{N}$, so the logarithm is well defined along the entire sequence and at the limit point. Therefore, by continuity of the logarithm on
\((0,\infty)\),
\[
\log\Phi_{\rho_n,\sigma_n}(\alpha-1)
\longrightarrow
\log\Phi_{\rho,\sigma}(\alpha-1).
\]
Consequently,
\[
S_\alpha^Q(\rho_n\|\sigma_n)
\longrightarrow
S_\alpha^Q(\rho\|\sigma).
\]
This completes the proof of continuity of $S^Q_\alpha(\rho\|\sigma)$
on $\mathcal{D}_+(\mathcal{H}) \times \mathcal{D}_+(\mathcal{H})$.
\end{proof}

\end{document}
